\documentclass[11pt]{article}

\usepackage{amsmath,amssymb,amsthm,mathtools}
\usepackage[a4paper,margin=1in]{geometry}
\usepackage[hidelinks]{hyperref}
\usepackage{enumitem}
\usepackage{xcolor}

\newtheorem{proposition}{Proposition}
\newtheorem{remark}{Remark}

\title{Feedback on the Revised Paper\\
\large Especially on the Section ``Aperiodicity for Irrational Slope''}
\author{}
\date{\today}

\begin{document}
\maketitle

\section*{Overall assessment}

The rational part of the paper remains in good shape.  The new figures are
helpful, the added converse statement in Step~2 is a real improvement, and
the clarifications concerning the \(\tilde p\)-coordinate make the set-valued
case easier to read.

The main issue is the new proposition on irrational slopes.  The statement is
plausible, and in the usual cut-and-project formulation it is essentially the
expected result.  However, the proof as currently written contains a serious
unjustified step.  In particular, the following assertion is too fast:
\begin{quote}
Since \(S'\) is the orthogonal projection of \(\mathbb Z^2\cap\text{strip}\)
onto the line, this translational symmetry must lift to a lattice translation
\((x,y)\mapsto(x+m,y+n)\) that preserves the strip.
\end{quote}
A period of the projected set gives a translation symmetry in physical space,
but it does not automatically follow, without proof, that this symmetry is
induced by one fixed lattice vector preserving the accepted set.  Fortunately,
for irrational slope this can be repaired cleanly using the injectivity of the
physical projection on \(\mathbb Z^2\) and the density of the internal
projection.

\section*{Main recommendation}

I recommend replacing the current proposition and proof by a more standard
cut-and-project argument.  The key points are:
\begin{enumerate}[label=\arabic*.]
    \item assume \(a>0\), not merely \(a\in\mathbb R\setminus\mathbb Q\);
    \item introduce physical and internal coordinates explicitly;
    \item use irrationality to prove injectivity of the physical projection;
    \item derive a genuine lattice translation from a projected period;
    \item use density of the internal projection to show that no nonzero
          internal shift can preserve the window interval.
\end{enumerate}

The following replacement is substantially safer.

\section*{Suggested replacement for the irrational-slope proposition}

\begin{proposition}[Aperiodicity for irrational slope]\label{prop:irrational_replacement}
Let \(a \in \mathbb{R}_{>0}\setminus\mathbb{Q}\) and \(\omega>0\).  Consider
the analogous strip projection with line \(f(x)=ax\) and strip
\[
    a(x-\omega)\le y\le ax+\omega.
\]
Then the projected gap sequence has no finite period.  Since the projection
map is injective on \(\mathbb Z^2\) for irrational \(a\), the multiset and set
conventions coincide in this case.
\end{proposition}

\begin{proof}
It is convenient to use the physical and internal coordinates
\[
    \pi(x,y) := x + ay,
    \qquad
    \pi^\star(x,y) := y - ax.
\]
The physical coordinate \(\pi\) determines the order of the orthogonal
projections onto the line \(y=ax\), up to multiplication by a positive
constant.  The strip condition is equivalent to
\[
    -a\omega \le \pi^\star(x,y) \le \omega.
\]
Thus the accepted lattice points are precisely
\[
    A := \{(x,y)\in\mathbb Z^2 : \pi^\star(x,y)\in W\},
    \qquad
    W := [-a\omega,\omega].
\]

First note that \(\pi\) is injective on \(\mathbb Z^2\).  Indeed, if
\(\pi(x_1,y_1)=\pi(x_2,y_2)\), then
\[
    x_1-x_2 = -a(y_1-y_2),
\]
which, since \(a\) is irrational and all entries are integers, forces
\(x_1=x_2\) and \(y_1=y_2\).

Suppose, for contradiction, that the projected gap sequence has a finite
period \(\lambda\).  Let
\[
    \Lambda := \pi(A)\subset\mathbb R
\]
be the projected set, and let \((s_i)_{i\in\mathbb Z}\) be its increasing
enumeration.  If the gap sequence has period \(\lambda\), then
\[
    s_{i+\lambda}-s_i = \sigma
\]
for all \(i\), where \(\sigma>0\) is the sum of one full period of consecutive
gaps.  Hence
\[
    \Lambda+\sigma=\Lambda.
\]

Choose any \(z_0\in A\).  Since \(\pi(z_0)+\sigma\in\Lambda\), there exists
\(z_1\in A\) such that
\[
    \pi(z_1)=\pi(z_0)+\sigma.
\]
Set \(v:=z_1-z_0\in\mathbb Z^2\).  Then \(\pi(v)=\sigma\).

We claim that \(A+v=A\).  Let \(z\in A\).  Since
\(\pi(z)+\sigma\in\Lambda\), there exists \(z'\in A\) with
\[
    \pi(z')=\pi(z)+\sigma=\pi(z+v).
\]
By injectivity of \(\pi\) on \(\mathbb Z^2\), this gives \(z'=z+v\).  Hence
\(z+v\in A\).  Applying the same argument to the period \(-\sigma\) gives
the reverse inclusion, so \(A+v=A\).

Write \(v=(m,n)\).  In internal coordinates, translation by \(v\) changes
\(\pi^\star\) by
\[
    \tau := \pi^\star(v)=n-am.
\]
Thus
\[
    \bigl(W\cap \pi^\star(\mathbb Z^2)\bigr)+\tau
    =
    W\cap \pi^\star(\mathbb Z^2).
\]
But
\[
    \pi^\star(\mathbb Z^2)=\{n-am:m,n\in\mathbb Z\}
\]
is dense in \(\mathbb R\), because \(a\) is irrational.

We now show that \(\tau=0\).  If \(\tau>0\), then by density we can choose
\(t\in\pi^\star(\mathbb Z^2)\cap W\) arbitrarily close to the upper endpoint
\(\omega\), with \(t+\tau>\omega\).  This contradicts the invariance above.
Similarly, if \(\tau<0\), choose
\(t\in\pi^\star(\mathbb Z^2)\cap W\) arbitrarily close to the lower endpoint
\(-a\omega\), with \(t+\tau<-a\omega\).  This is again a contradiction.
Hence \(\tau=0\).

Therefore
\[
    n-am=0.
\]
Since \(a\) is irrational and \(m,n\in\mathbb Z\), this forces \(m=n=0\).
Hence \(v=0\), and consequently
\[
    \sigma=\pi(v)=0,
\]
contradicting \(\sigma>0\).  Thus the projected gap sequence has no finite
period.
\end{proof}

\section*{Comments on the replacement proof}

This proof makes explicit the step that was implicit in the original draft.
The projected period first gives a physical translation by \(\sigma\).  Because
\(\pi\) is injective on the lattice for irrational \(a\), this physical
translation corresponds to one fixed lattice vector \(v\).  Once the lift to
\(v\) is justified, the internal coordinate shows that the acceptance window
would have to be invariant under a shift \(\tau=n-am\).  Since the internal
projection of the lattice is dense and the window is a compact interval with
endpoints, the only possible such shift is \(\tau=0\), which would force
\(a=n/m\) unless \(v=0\).  This contradicts the existence of a positive period.

\section*{Suggested revision of the dichotomy remark}

The current remark says, in effect,
\[
    \text{finite period}
    \quad\Longleftrightarrow\quad
    \text{rational slope}.
\]
This is fine in spirit, but I would phrase it slightly more carefully, because
the body of the paper is formulated mainly for positive rational slopes
\(\alpha/\beta\).  A safer version is:

\begin{quote}
Within the family of positive slopes considered here, the rational and
irrational cases therefore exhibit a sharp dichotomy: rational slopes give
periodic projected gap sequences, while irrational slopes give no finite
period.
\end{quote}

If you keep the equivalence notation, add explicitly that it is meant for the
present strip model with positive slopes.

\section*{Suggested revision in the introduction}

The introduction currently says that irrational slopes are ``typically
aperiodic''.  After the repaired proposition, you can state this more directly
for the present model.  I suggest:

\begin{quote}
When the slope is irrational, the resulting objects fall into the classical
aperiodic cut-and-project setting; for the present strip model we record the
corresponding absence of a finite gap period in
Proposition~\ref{prop:irrational}.
\end{quote}

This avoids the vague word ``typically''.

\section*{Placement recommendation}

The irrational-slope proposition is not really a corollary of the rational
period formula.  It is a separate statement about the complementary regime.
Thus I would either rename the section
\begin{quote}
\verb|\section{Corollaries}|
\end{quote}
to something like
\begin{quote}
\verb|\section{Consequences and the Irrational Case}|
\end{quote}
or, preferably, give the irrational statement its own short section:
\begin{quote}
\verb|\section{The Irrational Case}\label{sec:irrational}|
\end{quote}
This would make the structure clearer.

\section*{Other comments on the paper}

\begin{enumerate}[label=\arabic*.]
    \item The phrase ``strip around a line'' should still be changed to
    ``strip parallel to a line'', since the strip is not generally symmetric
    around the line \(y=(\alpha/\beta)x\).

    \item The statement ``the underlying set is \(\mathbb Z\)'' should always
    be qualified as ``in the \(\tilde p\)-coordinate''.  The current version
    mostly does this correctly.

    \item The figures are helpful, especially the track decomposition figure.
    For a specialist audience, however, you may consider reducing their size
    or keeping only the second figure.  This is a matter of presentation, not
    correctness.

    \item The empirical verification remarks are now appropriately framed as
    consistency checks rather than evidence needed for the proof.  This is
    good.

    \item The Lean table remains impressive but somewhat fragile if it contains
    exact line numbers.  If you keep the line numbers, it would be best to cite
    a repository commit hash.  Otherwise, consider omitting the line-number
    column.

    \item The acknowledgment should perhaps avoid the phrase ``proof
    construction''.  A more professional version would be:
    \begin{quote}
    The author used AI assistants during exploratory computations, drafting,
    and the preparation of the Lean formalization.  The author has
    independently verified the mathematical content and accepts full
    responsibility for the results and any errors.
    \end{quote}
\end{enumerate}

\section*{Final verdict}

The rational part of the paper is substantially improved and appears coherent.
The new irrational-slope section is a valuable addition, because it gives a
clear conceptual boundary between the rational periodic regime and the
irrational aperiodic regime.  However, the original proof of the irrational
proposition should not be sent in its present form.  It should be replaced by
an argument based on physical injectivity, internal density, and the absence of
nontrivial translations preserving an interval window, as suggested above.

\end{document}
